% Page content from a radiation safety manual in Spanish

\section*{Exposición controlada de una fuente radiactiva}

En esta sección, describimos el procedimiento controlado para la exposición de una fuente radiactiva utilizando un sistema de proyectores de radiografía. La ilustración muestra las etapas esenciales para exponer y retractar la fuente radiactiva de manera segura.

\begin{figure}[ht]
    \centering
    % \includegraphics[width=0.8\textwidth]{path_to_the_image}
    % TODO
    \caption{Etapas en la exposición y retracción de la fuente radiográfica}
    \label{fig:source_exposure}
\end{figure}

\begin{enumerate}
    \item \textbf{Posición almacenada:} El cable de control está retraído completamente permitiendo que la fuente de radiación se mantenga segura dentro del proyector. 
    \item \textbf{Fuente en tránsito:} Al girar la manivela, el cable de control desplaza la fuente desde el proyector hasta la posición de operación.
    \item \textbf{Fuente en el sitio radiográfico:} La fuente se sitúa en la posición focal, adecuada para la exposición radiográfica del objeto en inspección.
    \item \textbf{Retracción de la fuente:} Después de la exposición, la fuente se retracta de nuevo al proyector para asegurar su encapsulamiento y proteger a los operadores y al ambiente.
\end{enumerate}

\subsection*{Precauciones de seguridad}
Es crucial seguir todas las normativas de seguridad establecidas y verificar constantemente el funcionamiento efectivo de los mecanismos de control del proyector para prevenir cualquier exposición accidental o escape de la fuente radiactiva.

% Replace "path_to_the_image" with the actual path to the image when integrating this into the LaTeX document.

\newpage
